\documentclass[
  doc,
  floatsintext,
  longtable,
  a4paper,
  nolmodern,
  notxfonts,
  notimes,
  12pt,
  colorlinks=true,linkcolor=blue,citecolor=blue,urlcolor=blue]{apa7}

\usepackage{amsmath}
\usepackage{amssymb}

\geometry{inner=1in, outer=1in}
\fancyhfoffset[LE,RO]{0cm}


\usepackage[bidi=default]{babel}
\babelprovide[main,import]{spanish}


% get rid of language-specific shorthands (see #6817):
\let\LanguageShortHands\languageshorthands
\def\languageshorthands#1{}

\RequirePackage{longtable}
\RequirePackage{threeparttablex}

\makeatletter
\renewcommand{\paragraph}{\@startsection{paragraph}{4}{\parindent}%
	{0\baselineskip \@plus 0.2ex \@minus 0.2ex}%
	{-.5em}%
	{\normalfont\normalsize\bfseries\typesectitle}}

\renewcommand{\subparagraph}[1]{\@startsection{subparagraph}{5}{0.5em}%
	{0\baselineskip \@plus 0.2ex \@minus 0.2ex}%
	{-\z@\relax}%
	{\normalfont\normalsize\bfseries\itshape\hspace{\parindent}{#1}\textit{\addperi}}{\relax}}
\makeatother




\usepackage{longtable, booktabs, multirow, multicol, colortbl, hhline, caption, array, float, xpatch}
\usepackage{subcaption}


\renewcommand\thesubfigure{\Alph{subfigure}}
\setcounter{topnumber}{2}
\setcounter{bottomnumber}{2}
\setcounter{totalnumber}{4}
\renewcommand{\topfraction}{0.85}
\renewcommand{\bottomfraction}{0.85}
\renewcommand{\textfraction}{0.15}
\renewcommand{\floatpagefraction}{0.7}

\usepackage{tcolorbox}
\tcbuselibrary{listings,theorems, breakable, skins}
\usepackage{fontawesome5}

\definecolor{quarto-callout-color}{HTML}{909090}
\definecolor{quarto-callout-note-color}{HTML}{0758E5}
\definecolor{quarto-callout-important-color}{HTML}{CC1914}
\definecolor{quarto-callout-warning-color}{HTML}{EB9113}
\definecolor{quarto-callout-tip-color}{HTML}{00A047}
\definecolor{quarto-callout-caution-color}{HTML}{FC5300}
\definecolor{quarto-callout-color-frame}{HTML}{ACACAC}
\definecolor{quarto-callout-note-color-frame}{HTML}{4582EC}
\definecolor{quarto-callout-important-color-frame}{HTML}{D9534F}
\definecolor{quarto-callout-warning-color-frame}{HTML}{F0AD4E}
\definecolor{quarto-callout-tip-color-frame}{HTML}{02B875}
\definecolor{quarto-callout-caution-color-frame}{HTML}{FD7E14}

%\newlength\Oldarrayrulewidth
%\newlength\Oldtabcolsep


\usepackage{hyperref}



\usepackage{color}
\usepackage{fancyvrb}
\newcommand{\VerbBar}{|}
\newcommand{\VERB}{\Verb[commandchars=\\\{\}]}
\DefineVerbatimEnvironment{Highlighting}{Verbatim}{commandchars=\\\{\}}
% Add ',fontsize=\small' for more characters per line
\usepackage{framed}
\definecolor{shadecolor}{RGB}{241,243,245}
\newenvironment{Shaded}{\begin{snugshade}}{\end{snugshade}}
\newcommand{\AlertTok}[1]{\textcolor[rgb]{0.68,0.00,0.00}{#1}}
\newcommand{\AnnotationTok}[1]{\textcolor[rgb]{0.37,0.37,0.37}{#1}}
\newcommand{\AttributeTok}[1]{\textcolor[rgb]{0.40,0.45,0.13}{#1}}
\newcommand{\BaseNTok}[1]{\textcolor[rgb]{0.68,0.00,0.00}{#1}}
\newcommand{\BuiltInTok}[1]{\textcolor[rgb]{0.00,0.23,0.31}{#1}}
\newcommand{\CharTok}[1]{\textcolor[rgb]{0.13,0.47,0.30}{#1}}
\newcommand{\CommentTok}[1]{\textcolor[rgb]{0.37,0.37,0.37}{#1}}
\newcommand{\CommentVarTok}[1]{\textcolor[rgb]{0.37,0.37,0.37}{\textit{#1}}}
\newcommand{\ConstantTok}[1]{\textcolor[rgb]{0.56,0.35,0.01}{#1}}
\newcommand{\ControlFlowTok}[1]{\textcolor[rgb]{0.00,0.23,0.31}{\textbf{#1}}}
\newcommand{\DataTypeTok}[1]{\textcolor[rgb]{0.68,0.00,0.00}{#1}}
\newcommand{\DecValTok}[1]{\textcolor[rgb]{0.68,0.00,0.00}{#1}}
\newcommand{\DocumentationTok}[1]{\textcolor[rgb]{0.37,0.37,0.37}{\textit{#1}}}
\newcommand{\ErrorTok}[1]{\textcolor[rgb]{0.68,0.00,0.00}{#1}}
\newcommand{\ExtensionTok}[1]{\textcolor[rgb]{0.00,0.23,0.31}{#1}}
\newcommand{\FloatTok}[1]{\textcolor[rgb]{0.68,0.00,0.00}{#1}}
\newcommand{\FunctionTok}[1]{\textcolor[rgb]{0.28,0.35,0.67}{#1}}
\newcommand{\ImportTok}[1]{\textcolor[rgb]{0.00,0.46,0.62}{#1}}
\newcommand{\InformationTok}[1]{\textcolor[rgb]{0.37,0.37,0.37}{#1}}
\newcommand{\KeywordTok}[1]{\textcolor[rgb]{0.00,0.23,0.31}{\textbf{#1}}}
\newcommand{\NormalTok}[1]{\textcolor[rgb]{0.00,0.23,0.31}{#1}}
\newcommand{\OperatorTok}[1]{\textcolor[rgb]{0.37,0.37,0.37}{#1}}
\newcommand{\OtherTok}[1]{\textcolor[rgb]{0.00,0.23,0.31}{#1}}
\newcommand{\PreprocessorTok}[1]{\textcolor[rgb]{0.68,0.00,0.00}{#1}}
\newcommand{\RegionMarkerTok}[1]{\textcolor[rgb]{0.00,0.23,0.31}{#1}}
\newcommand{\SpecialCharTok}[1]{\textcolor[rgb]{0.37,0.37,0.37}{#1}}
\newcommand{\SpecialStringTok}[1]{\textcolor[rgb]{0.13,0.47,0.30}{#1}}
\newcommand{\StringTok}[1]{\textcolor[rgb]{0.13,0.47,0.30}{#1}}
\newcommand{\VariableTok}[1]{\textcolor[rgb]{0.07,0.07,0.07}{#1}}
\newcommand{\VerbatimStringTok}[1]{\textcolor[rgb]{0.13,0.47,0.30}{#1}}
\newcommand{\WarningTok}[1]{\textcolor[rgb]{0.37,0.37,0.37}{\textit{#1}}}

\providecommand{\tightlist}{%
  \setlength{\itemsep}{0pt}\setlength{\parskip}{0pt}}
\usepackage{longtable,booktabs,array}
\newcounter{none} % for unnumbered tables
\usepackage{calc} % for calculating minipage widths
% Correct order of tables after \paragraph or \subparagraph
\usepackage{etoolbox}
\makeatletter
\patchcmd\longtable{\par}{\if@noskipsec\mbox{}\fi\par}{}{}
\makeatother
% Allow footnotes in longtable head/foot
\IfFileExists{footnotehyper.sty}{\usepackage{footnotehyper}}{\usepackage{footnote}}
\makesavenoteenv{longtable}

\usepackage{graphicx}
\makeatletter
\newsavebox\pandoc@box
\newcommand*\pandocbounded[1]{% scales image to fit in text height/width
  \sbox\pandoc@box{#1}%
  \Gscale@div\@tempa{\textheight}{\dimexpr\ht\pandoc@box+\dp\pandoc@box\relax}%
  \Gscale@div\@tempb{\linewidth}{\wd\pandoc@box}%
  \ifdim\@tempb\p@<\@tempa\p@\let\@tempa\@tempb\fi% select the smaller of both
  \ifdim\@tempa\p@<\p@\scalebox{\@tempa}{\usebox\pandoc@box}%
  \else\usebox{\pandoc@box}%
  \fi%
}
% Set default figure placement to htbp
\def\fps@figure{htbp}
\makeatother







\usepackage{newtx}

\defaultfontfeatures{Scale=MatchLowercase}
\defaultfontfeatures[\rmfamily]{Ligatures=TeX,Scale=1}





\title{Salida de números con print()}


\shorttitle{SALIDA DE NÚMEROS CON PRINT()}


\usepackage{etoolbox}



\ccoppy{\textcopyright~2026}



\author{Edison Achalma}



\affiliation{
{Escuela Profesional de Economía, Universidad Nacional de San Cristóbal
de Huamanga}}




\leftheader{Achalma}

\date{2026-10-20}


\abstract{Este abstract será actualizado una vez que se complete el
contenido final del artículo. }



\authornote{\par{\addORCIDlink{Edison Achalma}{0000-0001-6996-3364}} 
\par{ }
\par{   El autor no tiene conflictos de interés que revelar.    Los
roles de autor se clasificaron utilizando la taxonomía de roles de
colaborador (CRediT; https://credit.niso.org/) de la siguiente
manera:  Edison Achalma:   conceptualización, metodología, análisis
formal, investigación, recursos, curación de
datos, redacción, visualización, supervisión, administración del
proyecto}
\par{La correspondencia relativa a este artículo debe dirigirse a Edison
Achalma, Escuela Profesional de Economía, Universidad Nacional de San
Cristóbal de Huamanga, Portal Independencia N°
57, Ayacucho, AYA 5001, Perú, Email: \href{mailto:elmer.achalma.09@unsch.edu.pe}{elmer.achalma.09@unsch.edu.pe}}
}

\makeatletter
\let\endoldlt\endlongtable
\def\endlongtable{
\hline
\endoldlt
}
\makeatother

\urlstyle{same}



\hypersetup{
  pdftitle={Salida de n\'{u}meros con print()},
  pdfauthor={Autor}
}
\makeatletter
\@ifpackageloaded{caption}{}{\usepackage{caption}}
\AtBeginDocument{%
\ifdefined\contentsname
  \renewcommand*\contentsname{Tabla de contenidos}
\else
  \newcommand\contentsname{Tabla de contenidos}
\fi
\ifdefined\listfigurename
  \renewcommand*\listfigurename{Lista de Figuras}
\else
  \newcommand\listfigurename{Lista de Figuras}
\fi
\ifdefined\listtablename
  \renewcommand*\listtablename{Lista de Tablas}
\else
  \newcommand\listtablename{Lista de Tablas}
\fi
\ifdefined\figurename
  \renewcommand*\figurename{Figura}
\else
  \newcommand\figurename{Figura}
\fi
\ifdefined\tablename
  \renewcommand*\tablename{Tabla}
\else
  \newcommand\tablename{Tabla}
\fi
}
\@ifpackageloaded{float}{}{\usepackage{float}}
\floatstyle{ruled}
\@ifundefined{c@chapter}{\newfloat{codelisting}{h}{lop}}{\newfloat{codelisting}{h}{lop}[chapter]}
\floatname{codelisting}{Listado}
\newcommand*\listoflistings{\listof{codelisting}{Lista de Listados}}
\makeatother
\makeatletter
\makeatother
\makeatletter
\@ifpackageloaded{caption}{}{\usepackage{caption}}
\@ifpackageloaded{subcaption}{}{\usepackage{subcaption}}
\makeatother
\makeatletter
\@ifpackageloaded{fontawesome5}{}{\usepackage{fontawesome5}}
\makeatother

% From https://tex.stackexchange.com/a/645996/211326
%%% apa7 doesn't want to add appendix section titles in the toc
%%% let's make it do it
\makeatletter
\xpatchcmd{\appendix}
  {\par}
  {\addcontentsline{toc}{section}{\@currentlabelname}\par}
  {}{}
\makeatother

%% Disable longtable counter
%% https://tex.stackexchange.com/a/248395/211326

\usepackage{etoolbox}

\makeatletter
\patchcmd{\LT@caption}
  {\bgroup}
  {\bgroup\global\LTpatch@captiontrue}
  {}{}
\patchcmd{\longtable}
  {\par}
  {\par\global\LTpatch@captionfalse}
  {}{}
\apptocmd{\endlongtable}
  {\ifLTpatch@caption\else\addtocounter{table}{-1}\fi}
  {}{}
\newif\ifLTpatch@caption
\makeatother

\begin{document}

\maketitle


\hypertarget{toc}{}
\tableofcontents
\newpage
\section[Introduction]{Salida de números con print()}

\setcounter{secnumdepth}{3}

\setlength\LTleft{0pt}




\subsection{🎯 Motivación}\label{motivaciuxf3n}

Los números son el alma de la programación: cálculos, mediciones,
estadísticas, todo gira en torno a ellos. Saber imprimir números
correctamente te permite ver resultados de operaciones, depurar cálculos
y mostrar información numérica a los usuarios. La diferencia clave con
el texto es que \textbf{los números no llevan comillas} --- y entender
esa diferencia es fundamental.

\subsection{💡 Intuición}\label{intuiciuxf3n}

Piensa en \texttt{print()} como una calculadora que también habla. Si le
pides que imprima \texttt{3}, simplemente muestra el número 3. Si le
pides \texttt{3\ +\ 3}, primero hace la cuenta (como una calculadora) y
luego muestra el resultado. Pero si le pides \texttt{"3"} (con
comillas), lo trata como texto, como si fuera una palabra escrita con el
símbolo ``3''.

\subsection{🧠 Modelo mental}\label{modelo-mental}

Python distingue entre:

\begin{enumerate}
\def\labelenumi{\arabic{enumi}.}
\tightlist
\item
  \textbf{Números (int, float)}: valores matemáticos. \textbf{Van sin
  comillas.}
\item
  \textbf{Texto (str)}: caracteres. \textbf{Van con comillas.}
\end{enumerate}

Cuando usas \texttt{print()}:

\begin{itemize}
\tightlist
\item
  Si pasas un número, lo muestra como número.
\item
  Si pasas una expresión matemática, \textbf{la calcula primero} y
  muestra el resultado.
\item
  Si pasas texto con comillas, lo muestra literalmente.
\end{itemize}

\subsection{📐 Desarrollo}\label{desarrollo}

\subsubsection{1. Imprimir números
simples}\label{imprimir-nuxfameros-simples}

\begin{Shaded}
\begin{Highlighting}[]
\BuiltInTok{print}\NormalTok{(}\DecValTok{3}\NormalTok{)      }\CommentTok{\# Número entero}
\BuiltInTok{print}\NormalTok{(}\FloatTok{3.14}\NormalTok{)   }\CommentTok{\# Número decimal (float)}
\BuiltInTok{print}\NormalTok{(}\OperatorTok{{-}}\DecValTok{5}\NormalTok{)     }\CommentTok{\# Número negativo}
\CommentTok{\# Salida:}
\CommentTok{\# 3}
\CommentTok{\# 3.14}
\CommentTok{\# {-}5}
\end{Highlighting}
\end{Shaded}

\subsubsection{\texorpdfstring{2. Operaciones matemáticas dentro de
\texttt{print()}}{2. Operaciones matemáticas dentro de print()}}\label{operaciones-matemuxe1ticas-dentro-de-print}

Python evalúa las operaciones antes de imprimir:

\begin{Shaded}
\begin{Highlighting}[]
\BuiltInTok{print}\NormalTok{(}\DecValTok{5} \OperatorTok{+} \DecValTok{3}\NormalTok{)      }\CommentTok{\# Suma}
\BuiltInTok{print}\NormalTok{(}\DecValTok{10} \OperatorTok{{-}} \DecValTok{4}\NormalTok{)     }\CommentTok{\# Resta}
\BuiltInTok{print}\NormalTok{(}\DecValTok{6} \OperatorTok{*} \DecValTok{7}\NormalTok{)      }\CommentTok{\# Multiplicación}
\BuiltInTok{print}\NormalTok{(}\DecValTok{15} \OperatorTok{/} \DecValTok{3}\NormalTok{)     }\CommentTok{\# División (resultado decimal)}
\BuiltInTok{print}\NormalTok{(}\DecValTok{10} \OperatorTok{//} \DecValTok{3}\NormalTok{)    }\CommentTok{\# División entera (trunca decimales)}
\BuiltInTok{print}\NormalTok{(}\DecValTok{10} \OperatorTok{\%} \DecValTok{3}\NormalTok{)     }\CommentTok{\# Módulo (resto de la división)}
\BuiltInTok{print}\NormalTok{(}\DecValTok{2} \OperatorTok{**} \DecValTok{4}\NormalTok{)     }\CommentTok{\# Potencia}
\CommentTok{\# Salida:}
\CommentTok{\# 8}
\CommentTok{\# 6}
\CommentTok{\# 42}
\CommentTok{\# 5.0}
\CommentTok{\# 3}
\CommentTok{\# 1}
\CommentTok{\# 16}
\end{Highlighting}
\end{Shaded}

\subsubsection{3. Combinar números y
texto}\label{combinar-nuxfameros-y-texto}

Usa \textbf{comas} para separar texto y números en el mismo
\texttt{print()}:

\begin{Shaded}
\begin{Highlighting}[]
\BuiltInTok{print}\NormalTok{(}\StringTok{"Tengo"}\NormalTok{, }\DecValTok{25}\NormalTok{, }\StringTok{"años"}\NormalTok{)}
\CommentTok{\# Salida: Tengo 25 años}

\BuiltInTok{print}\NormalTok{(}\StringTok{"El resultado es"}\NormalTok{, }\DecValTok{10} \OperatorTok{+} \DecValTok{5}\NormalTok{)}
\CommentTok{\# Salida: El resultado es 15}

\BuiltInTok{print}\NormalTok{(}\StringTok{"Precio:"}\NormalTok{, }\DecValTok{99}\NormalTok{, }\StringTok{"€"}\NormalTok{)}
\CommentTok{\# Salida: Precio: 99 €}
\end{Highlighting}
\end{Shaded}

\subsubsection{4. Texto vs números: la diferencia
clave}\label{texto-vs-nuxfameros-la-diferencia-clave}

\begin{Shaded}
\begin{Highlighting}[]
\BuiltInTok{print}\NormalTok{(}\DecValTok{5}\NormalTok{)       }\CommentTok{\# Muestra el número 5}
\BuiltInTok{print}\NormalTok{(}\StringTok{"5"}\NormalTok{)     }\CommentTok{\# Muestra el carácter "5" (texto)}
\BuiltInTok{print}\NormalTok{(}\DecValTok{5} \OperatorTok{+} \DecValTok{3}\NormalTok{)   }\CommentTok{\# Calcula 5+3, muestra 8}
\BuiltInTok{print}\NormalTok{(}\StringTok{"5+3"}\NormalTok{)   }\CommentTok{\# Muestra literalmente "5+3"}
\CommentTok{\# Salida:}
\CommentTok{\# 5}
\CommentTok{\# 5}
\CommentTok{\# 8}
\CommentTok{\# 5+3}
\end{Highlighting}
\end{Shaded}

\subsubsection{\texorpdfstring{5. Múltiples números en un
\texttt{print()}}{5. Múltiples números en un print()}}\label{muxfaltiples-nuxfameros-en-un-print}

\begin{Shaded}
\begin{Highlighting}[]
\BuiltInTok{print}\NormalTok{(}\DecValTok{10}\NormalTok{, }\DecValTok{20}\NormalTok{, }\DecValTok{30}\NormalTok{)}
\CommentTok{\# Salida: 10 20 30}

\BuiltInTok{print}\NormalTok{(}\DecValTok{1}\NormalTok{, }\DecValTok{2}\NormalTok{, }\DecValTok{3}\NormalTok{, sep}\OperatorTok{=}\StringTok{" {-} "}\NormalTok{)}
\CommentTok{\# Salida: 1 {-} 2 {-} 3}

\BuiltInTok{print}\NormalTok{(}\DecValTok{5}\NormalTok{, }\DecValTok{10}\NormalTok{, }\DecValTok{15}\NormalTok{, sep}\OperatorTok{=}\StringTok{", "}\NormalTok{)}
\CommentTok{\# Salida: 5, 10, 15}
\end{Highlighting}
\end{Shaded}

\subsection{🔬 Ejemplos}\label{ejemplos}

\textbf{Mínimo ---} Números y operaciones básicas:

\begin{Shaded}
\begin{Highlighting}[]
\BuiltInTok{print}\NormalTok{(}\DecValTok{100}\NormalTok{)}
\BuiltInTok{print}\NormalTok{(}\DecValTok{50} \OperatorTok{+} \DecValTok{50}\NormalTok{)}
\BuiltInTok{print}\NormalTok{(}\DecValTok{10} \OperatorTok{*} \DecValTok{2}\NormalTok{)}
\CommentTok{\# Salida:}
\CommentTok{\# 100}
\CommentTok{\# 100}
\CommentTok{\# 20}
\end{Highlighting}
\end{Shaded}

\textbf{Progresivo ---} Cálculos más complejos:

\begin{Shaded}
\begin{Highlighting}[]
\BuiltInTok{print}\NormalTok{(}\DecValTok{15} \OperatorTok{+} \DecValTok{3} \OperatorTok{*} \DecValTok{4}\NormalTok{)        }\CommentTok{\# Prioridad: multiplicación primero}
\BuiltInTok{print}\NormalTok{((}\DecValTok{15} \OperatorTok{+} \DecValTok{3}\NormalTok{) }\OperatorTok{*} \DecValTok{4}\NormalTok{)      }\CommentTok{\# Paréntesis cambian el orden}
\CommentTok{\# Salida:}
\CommentTok{\# 27   (15 + 12)}
\CommentTok{\# 72   (18 * 4)}
\end{Highlighting}
\end{Shaded}

\textbf{Real ---} Calculando propina en un restaurante:

\begin{Shaded}
\begin{Highlighting}[]
\NormalTok{cuenta }\OperatorTok{=} \FloatTok{45.50}
\NormalTok{propina\_porcentaje }\OperatorTok{=} \DecValTok{15}
\NormalTok{propina }\OperatorTok{=}\NormalTok{ cuenta }\OperatorTok{*}\NormalTok{ propina\_porcentaje }\OperatorTok{/} \DecValTok{100}
\NormalTok{total }\OperatorTok{=}\NormalTok{ cuenta }\OperatorTok{+}\NormalTok{ propina}

\BuiltInTok{print}\NormalTok{(}\StringTok{"Cuenta:"}\NormalTok{, cuenta, }\StringTok{"€"}\NormalTok{)}
\BuiltInTok{print}\NormalTok{(}\StringTok{"Propina (15\%):"}\NormalTok{, propina, }\StringTok{"€"}\NormalTok{)}
\BuiltInTok{print}\NormalTok{(}\StringTok{"Total a pagar:"}\NormalTok{, total, }\StringTok{"€"}\NormalTok{)}
\CommentTok{\# Salida:}
\CommentTok{\# Cuenta: 45.5 €}
\CommentTok{\# Propina (15\%): 6.825 €}
\CommentTok{\# Total a pagar: 52.325 €}
\end{Highlighting}
\end{Shaded}

\subsection{⚠️ Errores comunes}\label{errores-comunes}

{\def\LTcaptype{none} % do not increment counter
\begin{longtable}[]{@{}
  >{\raggedright\arraybackslash}p{(\linewidth - 4\tabcolsep) * \real{0.1832}}
  >{\raggedright\arraybackslash}p{(\linewidth - 4\tabcolsep) * \real{0.4198}}
  >{\raggedright\arraybackslash}p{(\linewidth - 4\tabcolsep) * \real{0.3969}}@{}}
\toprule\noalign{}
\begin{minipage}[b]{\linewidth}\raggedright
Error
\end{minipage} & \begin{minipage}[b]{\linewidth}\raggedright
Síntoma
\end{minipage} & \begin{minipage}[b]{\linewidth}\raggedright
Cómo detectarlo
\end{minipage} \\
\midrule\noalign{}
\endhead
\bottomrule\noalign{}
\endlastfoot
\textbf{Número con comillas} & \texttt{print("5")} muestra texto ``5'',
no número & Si necesitas hacer operaciones, \textbf{no} uses comillas \\
\textbf{Texto sin comillas} & \texttt{print(Hola)} →
\texttt{NameError:\ name\ \textquotesingle{}Hola\textquotesingle{}\ is\ not\ defined}
& El texto siempre va entre comillas \\
\textbf{Dividir por cero} & \texttt{print(5\ /\ 0)} →
\texttt{ZeroDivisionError:\ division\ by\ zero} & Asegura que el
denominador no sea 0 \\
\textbf{Operaciones en texto} & \texttt{print("5+3")} muestra ``5+3'',
no 8 & Las comillas hacen que todo sea literal \\
\textbf{Confundir separador} & \texttt{print("Hola"\ 25)} →
\texttt{SyntaxError} & Usa comas para separar:
\texttt{print("Hola",\ 25)} \\
\end{longtable}
}

\subsection{✅ Buenas prácticas}\label{buenas-pruxe1cticas}

\begin{itemize}
\item
  \textbf{Números sin comillas}: siempre que quieras un valor numérico o
  una operación.
\item
  \textbf{Usa paréntesis} para clarificar el orden de las operaciones.
\item
  \textbf{Nombra variables} en lugar de números ``mágicos'':

\begin{Shaded}
\begin{Highlighting}[]
\CommentTok{\# ❌ Poco legible}
\BuiltInTok{print}\NormalTok{(}\FloatTok{45.50} \OperatorTok{*} \FloatTok{0.15}\NormalTok{)}

\CommentTok{\# ✅ Legible}
\NormalTok{cuenta }\OperatorTok{=} \FloatTok{45.50}
\NormalTok{propina }\OperatorTok{=}\NormalTok{ cuenta }\OperatorTok{*} \FloatTok{0.15}
\BuiltInTok{print}\NormalTok{(propina)}
\end{Highlighting}
\end{Shaded}
\item
  \textbf{Separa texto y números con comas} para que Python los muestre
  correctamente.
\item
  \textbf{Usa f-strings} para formateo avanzado (lo veremos más
  adelante).
\end{itemize}

\subsection{📝 Ejercicio}\label{ejercicio}

\textbf{Básico} --- Calculadora de presupuesto

Escribe un programa que:

\begin{enumerate}
\def\labelenumi{\arabic{enumi}.}
\tightlist
\item
  Tenga un presupuesto de \textbf{500} euros
\item
  Un gasto en comida de \textbf{125.50} euros
\item
  Un gasto en transporte de \textbf{45.75} euros
\item
  Calcule el \textbf{total gastado} y el \textbf{saldo restante}
\item
  Imprima los resultados con mensajes descriptivos
\end{enumerate}

\textbf{Requisitos:} los números deben ser calculados, no escritos
directamente en los \texttt{print()}.

\begin{center}\rule{0.5\linewidth}{0.5pt}\end{center}

\textbf{Solución (no mires hasta haberlo intentado):}

\begin{Shaded}
\begin{Highlighting}[]
\NormalTok{presupuesto }\OperatorTok{=} \DecValTok{500}
\NormalTok{gasto\_comida }\OperatorTok{=} \FloatTok{125.50}
\NormalTok{gasto\_transporte }\OperatorTok{=} \FloatTok{45.75}

\NormalTok{total\_gastado }\OperatorTok{=}\NormalTok{ gasto\_comida }\OperatorTok{+}\NormalTok{ gasto\_transporte}
\NormalTok{saldo }\OperatorTok{=}\NormalTok{ presupuesto }\OperatorTok{{-}}\NormalTok{ total\_gastado}

\BuiltInTok{print}\NormalTok{(}\StringTok{"Presupuesto inicial:"}\NormalTok{, presupuesto, }\StringTok{"€"}\NormalTok{)}
\BuiltInTok{print}\NormalTok{(}\StringTok{"Gasto en comida:"}\NormalTok{, gasto\_comida, }\StringTok{"€"}\NormalTok{)}
\BuiltInTok{print}\NormalTok{(}\StringTok{"Gasto en transporte:"}\NormalTok{, gasto\_transporte, }\StringTok{"€"}\NormalTok{)}
\BuiltInTok{print}\NormalTok{(}\StringTok{"Total gastado:"}\NormalTok{, total\_gastado, }\StringTok{"€"}\NormalTok{)}
\BuiltInTok{print}\NormalTok{(}\StringTok{"Saldo restante:"}\NormalTok{, saldo, }\StringTok{"€"}\NormalTok{)}
\end{Highlighting}
\end{Shaded}

\textbf{Salida esperada:}

\begin{verbatim}
Presupuesto inicial: 500 €
Gasto en comida: 125.5 €
Gasto en transporte: 45.75 €
Total gastado: 171.25 €
Saldo restante: 328.75 €
\end{verbatim}

\subsection{🗂️ Síntesis}\label{suxedntesis}

\begin{itemize}
\tightlist
\item
  \textbf{Números sin comillas}: \texttt{print(5)} muestra el número 5.
\item
  \textbf{Texto con comillas}: \texttt{print("5")} muestra el texto
  ``5''.
\item
  \textbf{Operaciones matemáticas} se calculan antes de imprimir.
\item
  \textbf{Mezcla texto y números} con comas:
  \texttt{print("Edad:",\ 30)}.
\item
  \textbf{Jerarquía de operadores}: respeta la matemática estándar
  (paréntesis primero).
\item
  \textbf{Puedes usar \texttt{sep} y \texttt{end}} para personalizar el
  formato de salida.
\end{itemize}

\subsection{🔗 Qué sigue}\label{quuxe9-sigue}

\begin{itemize}
\tightlist
\item
  \textbf{Relacionado:} variables (guardar números en variables), tipos
  de datos (int, float, str)
\item
  \textbf{Siguiente tema:} Variables y tipos de datos --- aprenderás a
  almacenar y manipular números y texto de forma estructurada.
\end{itemize}

\section{Publicaciones Similares}\label{publicaciones-similares}

Si te interesó este artículo, te recomendamos que explores otros blogs y
recursos relacionados que pueden ampliar tus conocimientos. Aquí te dejo
algunas sugerencias:

\begin{enumerate}
\def\labelenumi{\arabic{enumi}.}
\tightlist
\item
  \href{https://numerus-scriptum.netlify.app/python/2020-06-19-instalacion-de-anaconda/index.pdf}{\faIcon{file-pdf}}
  \href{https://numerus-scriptum.netlify.app/python/2020-06-19-instalacion-de-anaconda}{Instalacion
  De Anaconda}
\item
  \href{https://numerus-scriptum.netlify.app/python/2020-06-20-configurar-entorno-virtual-python-anaconda/index.pdf}{\faIcon{file-pdf}}
  \href{https://numerus-scriptum.netlify.app/python/2020-06-20-configurar-entorno-virtual-python-anaconda}{Configurar
  Entorno Virtual Python Anaconda}
\item
  \href{https://numerus-scriptum.netlify.app/python/2021-04-17-01-introducion-a-la-programacion-con-python/index.pdf}{\faIcon{file-pdf}}
  \href{https://numerus-scriptum.netlify.app/python/2021-04-17-01-introducion-a-la-programacion-con-python}{01
  Introducion A La Programacion Con Python}
\item
  \href{https://numerus-scriptum.netlify.app/python/2021-05-31-02-variables-expresiones-y-statements-con-python/index.pdf}{\faIcon{file-pdf}}
  \href{https://numerus-scriptum.netlify.app/python/2021-05-31-02-variables-expresiones-y-statements-con-python}{02
  Variables Expresiones Y Statements Con Python}
\item
  \href{https://numerus-scriptum.netlify.app/python/2021-06-07-03-objetos-de-python/index.pdf}{\faIcon{file-pdf}}
  \href{https://numerus-scriptum.netlify.app/python/2021-06-07-03-objetos-de-python}{03
  Objetos De Python}
\item
  \href{https://numerus-scriptum.netlify.app/python/2021-06-14-04-ejecucion-condicional-con-python/index.pdf}{\faIcon{file-pdf}}
  \href{https://numerus-scriptum.netlify.app/python/2021-06-14-04-ejecucion-condicional-con-python}{04
  Ejecucion Condicional Con Python}
\item
  \href{https://numerus-scriptum.netlify.app/python/2021-06-21-05-iteraciones-con-python/index.pdf}{\faIcon{file-pdf}}
  \href{https://numerus-scriptum.netlify.app/python/2021-06-21-05-iteraciones-con-python}{05
  Iteraciones Con Python}
\item
  \href{https://numerus-scriptum.netlify.app/python/2021-08-16-06-funciones-con-python/index.pdf}{\faIcon{file-pdf}}
  \href{https://numerus-scriptum.netlify.app/python/2021-08-16-06-funciones-con-python}{06
  Funciones Con Python}
\item
  \href{https://numerus-scriptum.netlify.app/python/2021-08-23-07-dataframes-con-python/index.pdf}{\faIcon{file-pdf}}
  \href{https://numerus-scriptum.netlify.app/python/2021-08-23-07-dataframes-con-python}{07
  Dataframes Con Python}
\item
  \href{https://numerus-scriptum.netlify.app/python/2021-11-29-08-prediccion-y-metrica-de-performance-con-python/index.pdf}{\faIcon{file-pdf}}
  \href{https://numerus-scriptum.netlify.app/python/2021-11-29-08-prediccion-y-metrica-de-performance-con-python}{08
  Prediccion Y Metrica De Performance Con Python}
\item
  \href{https://numerus-scriptum.netlify.app/python/2021-12-06-09-metodos-de-machine-learning-para-clasificacion-con-python/index.pdf}{\faIcon{file-pdf}}
  \href{https://numerus-scriptum.netlify.app/python/2021-12-06-09-metodos-de-machine-learning-para-clasificacion-con-python}{09
  Metodos De Machine Learning Para Clasificacion Con Python}
\item
  \href{https://numerus-scriptum.netlify.app/python/2021-12-13-10-metodos-de-machine-learning-para-regresion-con-python/index.pdf}{\faIcon{file-pdf}}
  \href{https://numerus-scriptum.netlify.app/python/2021-12-13-10-metodos-de-machine-learning-para-regresion-con-python}{10
  Metodos De Machine Learning Para Regresion Con Python}
\item
  \href{https://numerus-scriptum.netlify.app/python/2025-05-10-visualizacion-de-datos-con-python/index.pdf}{\faIcon{file-pdf}}
  \href{https://numerus-scriptum.netlify.app/python/2025-05-10-visualizacion-de-datos-con-python}{Visualizacion
  De Datos Con Python}
\end{enumerate}

Esperamos que encuentres estas publicaciones igualmente interesantes y
útiles. ¡Disfruta de la lectura!






\end{document}
